\documentclass[11pt,a4paper]{article}
\usepackage[margin=24mm]{geometry}
\usepackage{amsmath,amssymb,amsthm,mathtools,array,booktabs}
\usepackage[hidelinks]{hyperref}
\newtheorem{theorem}{Theorem}[section]
\newtheorem{lemma}[theorem]{Lemma}
\newtheorem{corollary}[theorem]{Corollary}
\newtheorem{proposition}[theorem]{Proposition}
\theoremstyle{remark}
\newtheorem{remark}[theorem]{Remark}
\newcommand{\C}{\mathbb C}
\newcommand{\R}{\mathbb R}
\title{Mordell's distance product:\\published closure for six and seven points\\and new eight-point exclusions}
\author{}
\date{8 October 2026}
\begin{document}
\maketitle
\begin{abstract}
The unrestricted six- and seven-point values are complete using the
published physical classification of Faug\`ere--Svartz and fresh exact
retained-chart comparisons. Eight points remain open. New results
exclude every centrally symmetric physical stationary configuration
on two orthogonal axes, classify the previously untreated unbalanced
axis branch, and exclude three entire boundary subcharts of the
first-nonzero-fourth-coefficient chart. A further exact reduction
determines the coupled phases of every nonreflected three-radius
population \(1+3+4\) candidate in the global radius range from two
radii; its degenerate triangle branch is proved reflected.
The radii satisfy a further explicit degree-six algebraic curve.
This reduction does not exhaust that population.
The general even chart and unrestricted eight-point comparison are
not claimed complete. The earlier nonclassification regional results
remain valid and available in the companion addendum.
\end{abstract}
\tableofcontents

\section*{Normalizations and explicit earlier inputs}
For \(\sum_i|z_i|^2=n\), write
\(\Delta_n=\prod_{i<j}|z_i-z_j|^2\). The central polygon candidate is
\[
 C_n=\frac{n^{n(n-1)/2}}
 {(n-1)^{(n-1)(n-2)/2}},
 \qquad C_8=\frac{8^{28}}{7^{21}}.
\]
The three coordinate conventions used below are related by
\(u=\sqrt{7/8}\,z\), \(w=2u=\sqrt{7/2}\,z\) for eight points.
Their moments are respectively \(8,7,28\), and their force constants
are \(7/2,4,1\). Every change of convention is stated locally.

The global structural corollary in Section~2 uses the previously
proved no-origin two-level exclusion, global radius and subset-spread
bounds, and the completed centered three-level \(2+2+4\) certificate.
The corresponding source proofs are included as prerequisites in
the replay package. These are explicit inputs to that corollary.
The new axis instability theorem itself needs none of those global
cuts. Section~1 explicitly identifies the published classification
inputs; local root disks are not asserted to reconstruct them.
% Fragment: requires amsmath, amssymb, booktabs, theorem/proof environments,
% and the bibliography entries FS, FSdata, FSdataSeven.
\section{Six and seven points: closure using the published classification}
\label{sec:classification-bridge67}

Let $N_n=n(n-1)/2$, $\Delta_n(z)=\prod_{i<j}|z_i-z_j|^2$, and
\[
 C_6=\frac{6^{15}}{5^{10}},\qquad C_7=\frac{7^{21}}{6^{15}}.
\]
This section uses the published physical classification of
Faug\`ere--Svartz~\cite{FS}. Its finite arithmetic postprocessing was
replayed afresh in \texttt{classification\_bridge/replay}; no numerical
screening decision is used in the proof.

\paragraph{The classification input.}
We use the following consequence of the exact computation reported in
\cite[Sections~5.3--5.7]{FS}, with the archived author data
\cite{FSdata,FSdataSeven}: every collision-free physical stationary shape
for $n=6,7$, modulo complex scaling and permutation, is represented in
the retained coefficient charts described below. The computation includes
unstable and collinear equilibria. Section~5.3 uses first-nonzero-coefficient
charts and the collision saturation $hD=1$; its Remark~6 explains the
positive real rescaling after coefficient normalization. Section~5.4
filters the actual pointwise conjugate equations. These sections occur on
proceedings p.~176, and Section~5.7 on p.~177.

For $n=6$, the retained parameter degrees are
$22,28,12,4,2,1,1,6,2,1,1$, totaling $80$.
The literal $\texttt{N=7}$ header in the archived six-point file is a
metadata error; its variables are $e_2,\ldots,e_6$. The eleven forward
substitution checks into the six-point necessary equations were also
replayed. Forward substitution does not establish chart completeness.

For $n=7$, \cite[Section~7, Figure~4, p.~178]{FS} is an essential part of
the input. The $e_2=1$ component has degree
$526=6+18+356+82+64$. The figure declares the degree-$356$ branch physically
empty. Its polynomial, coordinate maps and arithmetic exclusion certificate
are absent from the local package. The retained degrees are
$1,1,2,6,18,82,64$, totaling $174$. The degree-$18$ branch is retained and
independently excluded here. Thus $174$ is the number of retained parameter
roots, not the number of all complex algebraic candidates. The missing
degree-$356$ physical exclusion and original chart completeness are cited
published inputs. Theorem~5 of~\cite{FS} supplies the symmetric equations;
the subsequent computation supplies the classification used here.

\begin{lemma}[Scaling and the original conjugate equations]
\label{lem:cb67-reduction}
A global maximizer of $\Delta_n$ on $\sum_i|z_i|^2=n$ is collision-free,
centered, and physical stationary. If $w_i=\alpha z_i$, $\alpha\ne0$,
$Q=\sum_i|w_i|^2$ and $P(t)=\prod_i(t-w_i)$, then
\begin{equation}\label{eq:cb67-force}
 QP''(w_i)-2N_n\overline{w_i}P'(w_i)=0\quad(1\le i\le n),
\end{equation}
and the moment-$n$ value of its shape is
\begin{equation}\label{eq:cb67-value}
 \mathcal V_n(w)=\left(\frac nQ\right)^{N_n}
                    \prod_{i<j}|w_i-w_j|^2.
\end{equation}
\end{lemma}
\begin{proof}
Compactness gives a maximizer; a regular $n$-gon gives a positive value,
so its points are distinct. If $c=n^{-1}\sum_i z_i\ne0$, translation by
$-c$ preserves all differences and reduces the moment to $n-n|c|^2<n$.
Rescaling to moment $n$ strictly increases the positive objective. Hence
$c=0$.

The Lagrange equations for $\log\Delta_n$ on the smooth sphere are
$F_i(z)=\lambda\overline{z_i}$, where
$F_i(z)=\sum_{j\ne i}(z_i-z_j)^{-1}$ and $\lambda$ is real.
Pairwise summation gives $\sum_i z_iF_i(z)=N_n$, so
$\lambda=N_n/n>0$. For $w=\alpha z$,
$F_i(w)=\lambda|\alpha|^{-2}\overline{w_i}$.
Writing $\kappa=\lambda|\alpha|^{-2}$ gives $\kappa Q=N_n$.
Since $P'(w_i)\ne0$ and $F_i(w)=P''(w_i)/(2P'(w_i))$, this proves
\eqref{eq:cb67-force}. Scaling to moment $n$ proves
\eqref{eq:cb67-value}.

The centered polynomial is
$P(t)=t^n+\sum_{k=2}^n(-1)^ke_k t^{n-k}$.
Its first nonzero $e_k$ can be normalized to $1$ by a complex scale
$\alpha^k e_k=1$. Earlier zero coefficients remain zero, and
\eqref{eq:cb67-value} is unchanged. Moreover $u=\sqrt\kappa\,w$ satisfies
$F_i(u)=\overline{u_i}$ in the standard physical units, so the positive-force
classification applies. Coefficient normalization then returns precisely
the homogeneous charts described in Remark~6 of~\cite{FS}.
\end{proof}

\paragraph{The rational certificate and its exact tests.}
The polynomial conventions are
\[
 \begin{split}
 P_6(t)&=t^6+e_2t^4-e_3t^3+e_4t^2-e_5t+e_6,\\
 P_7(t)&=t^7+e_2t^5-e_3t^4+e_4t^3-e_5t^2+e_6t-e_7.
 \end{split}
\]
All coefficient charts are parsed as rational polynomials rather than
executed. Parameter disks have radii $10^{-70}$ for $n=6$ and $10^{-250}$
for $n=7$; point-root disks have radii $10^{-12}$ and $10^{-10}$,
respectively. For a Gaussian rational $a$ put
$\|a\|_1=|\Re a|+|\Im a|$ and
$\|a\|_\infty=\max(|\Re a|,|\Im a|)$, so
$\|a\|_\infty\le |a|\le\|a\|_1$.

For a nominal polynomial $p_0$ with coefficient errors $\epsilon_j$ and
$a_k=p_0^{(k)}(c)/k!$, the strict test
\[
 \|a_0\|_1+\sum_{k\ge2}\|a_k\|_1r^k
       +\sum_j\epsilon_j(\|c\|_1+r)^j
 <r\|a_1\|_\infty
\]
is Rouch\'e comparison with $a_1(t-c)$ on $|t-c|=r$. It certifies exactly
one root counting multiplicity. Disjoint disks and degree counts certify
every retained parameter root and every point root. For high-degree
seven-point parameter polynomials the same remainder is bounded by
$r^2\sup_{|t|\le B}|p_0''(t)|/2$, with the rationally computed integer
$B=\lceil\|c\|_1+r\rceil$.
Coefficient errors include parameter variation and rational rounding of
the coefficient centers.

For point centers $c_i$ with common radius $r$, set
$Q_0=\sum_i|c_i|^2$ and
$E_Q=\sum_i(2\|c_i\|_1r+r^2)$.
At one point let $v_j=p_0^{(j)}(c_i)$ and let $E_j$ bound
$|P^{(j)}(w_i)-v_j|$, with parameter and point uncertainty included.
The residual in \eqref{eq:cb67-force} differs from
$Q_0v_2-2N_n\overline{c_i}v_1$ by at most
\[
 E_Q\|v_2\|_1+(Q_0+E_Q)E_2
 +2N_n\{r\|v_1\|_1+(\|c_i\|_1+r)E_1\}.
\]
If its nominal $\|\cdot\|_\infty$ exceeds this quantity at any one point,
the entire candidate is physically excluded. In particular this tests the
actual conjugate of each corresponding root, not an unordered conjugate
set of roots. When no exclusion is used, $Q_0>E_Q$ and
\begin{equation}\label{eq:cb67-upper}
 U_n=\left(\frac n{Q_0-E_Q}\right)^{N_n}
 \prod_{i<j}\bigl(|c_i-c_j|^2+4r\|c_i-c_j\|_1+4r^2\bigr)
\end{equation}
is a rigorous rational upper bound for \eqref{eq:cb67-value}.

\begin{proposition}[Fresh exhaustive retained-chart replay]
\label{prop:cb67-replay}
For $n=6$ the $80$ parameter roots and $480$ point-root disks split into
$69$ force exclusions, $10$ strictly smaller candidates, and the exact
central pentagon. For $n=7$ the $174$ retained parameter roots and $1218$
point-root disks split into $155$ force exclusions, $18$ strictly smaller
candidates, and the exact central hexagon. Every nonattaining physical
stationary shape has moment-$n$ value at most $6^6$ or $7^7$, respectively.
\end{proposition}
\begin{proof}
The fresh full runs of \path{certify_fs6_disks.py} \texttt{--verify} and
\path{certify_fs7_disks.py} \texttt{--verify}, under
$\texttt{python -B -S -O}$, check every strict inequality just described.
Their decisions use integers and $\texttt{fractions.Fraction}$, with
explicit exceptions active under $\texttt{-O}$. The resulting branch
counts and largest upward-rounded integer bounds are:
\begin{center}\begin{tabular}{rrrrrl}
\toprule
$n$&Chart&Parameters&Excluded&Bounded&Largest bound or exact chart\\
\midrule
6&1&22&20&2&$3170$\\
6&2&28&28&0&---\\
6&3&12&12&0&---\\
6&4&4&2&2&$42454$\\
6&5&2&2&0&---\\
6&6&1&0&1&$18699$\\
6&7&1&0&1&$18699$\\
6&8&6&4&2&$22400$ (other bound $133$)\\
6&9&2&1&1&$7225$\\
6&10&1&0&0&$P=t(t^5-1)$, attaining\\
6&11&1&0&1&$6^6=46656$, exact\\
\midrule
7&1&1&0&1&$7^7=823543$, exact\\
7&2&1&0&0&$P=t(t^6+1)$, attaining\\
7&3&2&1&1&$40253$\\
7&4&6&4&2&$159666$\\
7&5&18&18&0&---\\
7&6&82&78&4&$268962$\\
7&7&64&54&10&$822629$\\
\bottomrule
\end{tabular}\end{center}
The independent aggregation script checks all degree counts, decision
partitions, displayed maxima, unique attaining charts and the exact
integer comparisons
\[
 6^9=10077696>9765625=5^{10},\qquad
 7^{14}=678223072849>470184984576=6^{15}.
\]
Consequently $6^6<C_6$ and $7^7<C_7$. The differences are exactly
\[
 C_6-6^6=\frac{14559984576}{9765625}>0,\qquad
 C_7-7^7=\frac{171328311330611239}{470184984576}>0.
\]
These are gaps for other physical stationary shapes; no uniform fixed gap
is asserted for arbitrary configurations approaching an attaining one.
The auxiliary eleven-component substitution replay proves forward
containment only. Published completeness remains the stated input.
\end{proof}

\begin{theorem}[Complete six- and seven-point global values]
\label{thm:cb67-global}
Using the stated published physical classification, every configuration
with $\sum_i|z_i|^2=n$ satisfies
\[
 \Delta_6\le\frac{6^{15}}{5^{10}},\qquad
 \Delta_7\le\frac{7^{21}}{6^{15}}.
\]
Equality holds precisely for the origin together with a regular pentagon
of radius $\sqrt{6/5}$, or a regular hexagon of radius $\sqrt{7/6}$,
respectively, up to rotation and permutation.
\end{theorem}
\begin{proof}
A global maximizer satisfies Lemma~\ref{lem:cb67-reduction}.
Published physical coverage puts it in a retained chart, and
Proposition~\ref{prop:cb67-replay} bounds its value by $C_n$.
This also closes the previously untreated configurations with no origin:
if any such configuration exceeded $C_n$, the compact sphere would have
a global maximizer exceeding $C_n$, contrary to the same argument.

Attainment is checked directly, without inferring physical sufficiency
from a failed exclusion test. Put $m=n-1$ and take the roots of
$P(t)=t(t^m-a)$, $|a|=1$. Their moment is $m$.
Since $|P'(0)|=1$ and $|P'(w)|=m$ at each nonzero root,
the modulus of the monic discriminant is $m^m$, so their squared
pair-distance product is $m^m$.
At a nonzero root, $P''(w)/(2P'(w))=n/(2w)$, while the force at the origin
is zero. Thus the original physical equation also holds directly.
Scaling to moment $n$ gives
\[
 (n/m)^{N_n}m^m
 =\frac{n^{N_n}}{(n-1)^{(n-1)(n-2)/2}}=C_n\quad(n=6,7).
\]
Every equality configuration is itself a global maximizer. The strict
retained-chart comparisons force it into the unique attaining chart.
Moment normalization fixes the radius, and the remaining complex phase
is an overall rotation. This proves both equality classifications.
\end{proof}

% Section fragment. Requires amsmath, amssymb, amsthm and theorem environments.
% New complete branch; old proofs and generated resultants remain unchanged.
\section{Complete instability of centrally symmetric eight-point axis configurations}
\label{sec:eight-even-new-axes}

We use force-one coordinates $w$, so true stationarity is
\[
 \sum_{j\ne i}\frac1{w_i-w_j}=\overline{w_i},\qquad
 \sum_i|w_i|^2=28.
\]
For a centrally symmetric collision-free configuration with no origin,
write its eight roots as $\pm\sqrt{y_i}$, $1\le i\le4$, where
the $y_i$ are distinct and nonzero. Put
\[
 q(y)=\prod_{i=1}^4(y-y_i),\quad H=q'+2yq'',\quad T_i=|y_i|.
\]
The actual force equations, with their physical signs and pairing, are
\begin{equation}\label{eq:eight-even-new-physical}
 \frac{H(y_i)}{2q'(y_i)}=T_i
 =\frac12+2y_i\sum_{j\ne i}\frac1{y_i-y_j}.
\end{equation}
Pairwise summation gives $\sum_iT_i=14$. The moment-eight coordinates
are $z=\sqrt{2/7}\,w$, multiplying the product by $(2/7)^{28}$.
No unordered resultant-root matching is used in this section.

\subsection{The full symmetry-breaking Hessian and a universal test vector}
On perturbations assigning the same displacement $s_i$ to both
$\sqrt{y_i}$ and $-\sqrt{y_i}$, the inverse-square Hessian has block
\[
 C_{ij}=-\frac{2(y_i+y_j)}{(y_i-y_j)^2}\quad(i\ne j),\qquad
 C_{ii}=-\sum_{j\ne i}C_{ij}.
\]
Indeed, its off-diagonal entry is exactly
$-(\sqrt{y_i}-\sqrt{y_j})^{-2}-(\sqrt{y_i}+\sqrt{y_j})^{-2}$.
Thus $C=C^{\mathsf T}$ and $C\mathbf1=0$.
If $\sum_i s_i=0$, the eight-point displacement is centered;
its moment-tangent condition holds automatically because the
contributions of each opposite pair cancel. The second-order term
of $\log\Delta_8-\sum_i|w_i|^2$ on this displacement is
\[
 -2\|s\|^2-2\Re(s^{\mathsf T}Cs).
\]
Every phase multiple remains a legal tangent vector. Therefore an
actual local maximum must satisfy
\begin{equation}\label{eq:eight-even-new-C-contract}
 |s^{\mathsf T}Cs|\le\|s\|^2\quad(s\perp\mathbf1),\qquad
 \|C\|\le1.
\end{equation}
The equivalence follows from Takagi factorization, using $C\mathbf1=0$.
These perturbations break central symmetry and must be retained even
when the candidate polynomial is even.

The physical force identity gives
\[
 Cy=2T-7\mathbf1.
\]
When all four $y_i$ are real, define
\[
 m=\sum_i y_i,\quad S_2=\sum_i y_i^2=\sum_iT_i^2,\quad
 v=y-\frac m4\mathbf1.
\]
The centered input $v$ obeys
\begin{equation}\label{eq:eight-even-new-defect}
 \|v\|^2=S_2-m^2/4,\quad
 \|Cv\|^2=4S_2-196,\quad
 \|v\|^2-\|Cv\|^2=196-3S_2-m^2/4.
\end{equation}
The last quantity must be nonnegative at a full local maximum.

\begin{theorem}\label{thm:eight-even-new-axes-unstable}
Every collision-free, no-origin, centrally symmetric physical
stationary eight-point configuration supported on two orthogonal
lines through the origin is unstable in the full eight-point problem.
Equivalently, all physical stationary configurations with four real
quartic roots $y_i$ fail the necessary contraction condition
\eqref{eq:eight-even-new-C-contract}.
\end{theorem}
\begin{proof}
Rotate the two lines to the real and imaginary axes. Their squared
roots $y_i$ are real. Multiplying all $w_i$ by $\mathrm i$ exchanges
positive and negative $y_i$, preserving stationarity and instability.
It suffices to consider zero, one, or two positive quartic roots.

If all roots are positive, $m=14$ and $Cv=2v$; if all are negative,
$m=-14$ and $Cv=-2v$. Since the roots are distinct, $v\ne0$.
Thus $\|C\|\ge2>1$ in both cases.

Suppose there is one positive root $A$ and three negative roots.
Write
\[
 f(y)=y^3+by^2+cy+d,\qquad q(y)=(y-A)f(y).
\]
The physical moment gives $b=14-A$. At every negative root,
$H+2yq'=0$, so this polynomial has zero remainder modulo $f$.
After substitution of $b=14-A$, the three remainder coefficients,
in decreasing powers of $y$, are
\[
 -2(15A+2c-105),\quad
 2(3A^2+Ac-42A+3c-3d),\quad -Ac+4Ad+d.
\]
The first two imply
\[
 c=\frac{105-15A}{2},\qquad
 d=\frac{105-8A-3A^2}{2}.
\]
Already the first identity gives
\[
 m=2A-14,\qquad S_2=A^2+b^2-2c=2A^2-13A+91.
\]
Consequently
\[
 \|v\|^2=A^2+A+42,\qquad
 \|Cv\|^2=8A^2-52A+168,
\]
and the defect is the strictly negative exact square certificate
\begin{equation}\label{eq:eight-even-new-unbalanced-square}
 \|v\|^2-\|Cv\|^2
 =-7A^2+53A-126
 =-7\left(A-\frac{53}{14}\right)^2-\frac{719}{28}<0.
\end{equation}
This excludes the whole one-positive/three-negative physical branch,
not just a sample or a radial subinterval.

It remains to treat two positive and two negative roots. Order them
in the convex open chamber $y_1<y_2<0<y_3<y_4$. The physical
equations are exactly the gradient equations of
\[
 \Phi(y)=\sum_i\log|y_i|+
             4\sum_{i<j}\log|y_i-y_j|-2\sum_i|y_i|.
\]
Within this fixed-sign ordered chamber the absolute-value term is
affine, and for every nonzero real $h$,
\[
 h^{\mathsf T}(\nabla^2\Phi)h
 =-\sum_i\frac{h_i^2}{y_i^2}
   -4\sum_{i<j}\frac{(h_i-h_j)^2}{(y_i-y_j)^2}<0.
\]
Thus $\Phi$ is strictly concave, and there is at most one stationary
point in the chamber. No existence statement is needed here.
Negating and reversing the ordered roots preserves the chamber,
the objective, and the stationary equations. If an actual stationary
point exists, uniqueness forces its root set to be
$\{-R,-S,S,R\}$ with $R>S>0$.
The moment gives $R+S=7$, and the physical force at $R$ reads
\[
 R=\frac32+\frac{4R^2}{R^2-S^2}.
\]
With $D=R-S$ this is $5D^2=98$, so
$S_2=(R+S)^2+(R-S)^2=343/5$ and $m=0$. Equation
\eqref{eq:eight-even-new-defect} gives
\[
 \|v\|^2-\|Cv\|^2=-49/5<0.
\]
Equivalently $\|Cy\|^2/\|y\|^2=8/7>1$.
Every possible sign chamber is therefore unstable.
\end{proof}

\subsection{A complete physical classification of the unbalanced axis branch}
\begin{proposition}\label{prop:eight-even-new-axis-classification}
Up to rotation, permutation, and exchange of the two axes, there is
exactly one physical stationary configuration with two points on one
axis and six on the other. Its positive quartic root $A$ is the unique
positive root of
\[
 Q(A)=12A^3+20A^2-307A-105,
 \qquad 22/5<A<9/2,
\]
and its three negative roots are the roots of
\[
 f(y)=y^3+(14-A)y^2+
       \frac{105-15A}{2}y+\frac{105-8A-3A^2}{2}.
\]
They lie respectively in $(-9,-6)$, $(-6,-1)$, and $(-1,0)$.
The most negative root is below $-49/8$, also violating the global
force-one radius bound. The Hessian exclusion above needs no such bound.
\end{proposition}
\begin{proof}
The first two negative-root remainder coefficients force the displayed
$f$. Exact division and substitution then give both identities
\[
 (H+2yq')\bmod f=-Q(A)/2,\qquad
 H(A)-2Aq'(A)=-Q(A)/2.
\]
Thus every physical solution has $Q(A)=0$. Its coefficients have one
Descartes sign change, while $Q(22/5)<0<Q(9/2)$, giving exactly one
positive root.
On this rational interval the coefficients lie in
\[
 19/2<b<48/5,\quad 75/4<c<39/2,\quad 33/8<d<293/50.
\]
The monotonicity of $d=(105-8A-3A^2)/2$ is immediate for $A>0$.
Direct rational interval signs give
$f(-9)<0$, $f(-6)>0$, $f(-1)<0$, and $f(0)>0$ uniformly.
Therefore $f$ has three distinct negative real roots in the indicated
disjoint intervals and no others. Also $f(-49/8)>0$, placing its
leftmost root below $-49/8$.
These roots and $A$ are nonzero and distinct, hence $q'$ does not
vanish at them. The two exact force identities now give
$H/(2q')=-y=|y|$ on each negative root and $H/(2q')=A=|A|$ at
the positive root. For each square root $w$, this implies
$P''(w)/(2P'(w))=|w|^2/w=\overline w$, including both signs.
This proves existence, uniqueness and genuine physical pairing,
rather than only necessary squared or setwise equations.
\end{proof}

\subsection{A reflection exclusion specifically for global even candidates}
\begin{corollary}\label{cor:eight-even-new-no-reflection-global}
Every still-unexcluded no-origin centrally symmetric eight-point
global maximizer has four distinct nonzero radial levels and no
reflection symmetry axis.
\end{corollary}
\begin{proof}
The existing no-origin two-level theorem excludes all global candidates
with at most two radial levels. With three levels, central symmetry
centers each level automatically and forces the population $2+2+4$.
The completed three-level certificate
\texttt{n8\_algebra\_threelevel\_224.tex}, Proposition
\texttt{prop:eight-centered-threelevel-224}, excludes that complete
global branch. Its four-point-layer domain $3/8<c<49/32$ in
$u=\sqrt{7/8}\,z$ coordinates follows from the strict four-point
internal-moment bound and the global one-point radius bound.
This uses the completed $2+2+4$ certificate, not the earlier
structural reduction that left this case open.

Thus a remaining even global candidate must have four different
values of $|y_i|$. If it has a reflection axis, the centroid at zero
forces the axis through zero. Rotate it to the real axis. Its even
root polynomial $P$ and quartic $q$ then have real coefficients.
Any nonreal quartic root would have a distinct conjugate root with
the same modulus, impossible with four distinct moduli. All four
quartic roots must therefore be real, and Theorem
\ref{thm:eight-even-new-axes-unstable} rules out a local, hence a
global, maximum. This proves the reflection exclusion for global
candidates; it does not assert instability of every non-axis
reflection-symmetric stationary configuration.
\end{proof}

The program \path{axis_even_check.py} uses Python's standard library.
Its Fraction sparse-polynomial arithmetic passes
42 exact checks under \texttt{-B -S -O}. It verifies the original
force remainders, genuine sign matching, rational root intervals,
the complete square and balanced-sign defects, and independent
rational audits of the universal $C$ action.
The analytic strict-concavity and full Hessian implications are
proved above. The general complex quartic branch with four distinct
moduli and no reflection axis remains unresolved.

\subsection{Recalled two-square theorem used by the next exclusion}
The following proposition is an earlier proved input, reproduced here
for the first-fourth-coefficient argument. It is not a new result of
this addendum. Its block \(C\) is the same block defined above.
\begin{proposition}[All stationary two-square configurations are unstable]
\label{prop:n8-two-square-unstable}
Every collision-free stationary eight-point configuration consisting
of two concentric square rings is unstable under perturbations that
break central symmetry. In force-one scale the largest singular value
of the same-pair block $C$ has square
\[
\begin{aligned}
 \frac{64}{49}&\quad\hbox{(regular octagon)},\\
 \frac{532}{245}&\quad\hbox{(aligned unequal squares)},\\
 \frac87&\quad\hbox{(staggered unequal squares)}.
\end{aligned}
\]
In particular none is a local maximum among unrestricted eight-point
configurations.
\end{proposition}
\begin{proof}
Let $R,S>0$ be the squared ring radii. Rotate the first square so its
fourth powers are $R^2$. Multiplying its pointwise force equation by
the point gives
\[
 R=\frac32+\frac{4R^2}{R^2-X},
\]
where $|X|=S^2$ is the fourth power on the other square. Reality of
the left side forces $X$ real, hence $X=\varepsilon S^2$,
$\varepsilon=1$ or $-1$. The moment gives $R+S=7$.
For $\varepsilon=1$, collision-freeness requires $R\ne S$ and the
force equation gives $(R-S)^2=98/5$. For $\varepsilon=-1$, subtract
the force equations of the two rings to obtain either $R=S=7/2$ or
$R^2+S^2=28$, equivalently $(R-S)^2=7$.

The quartic roots are $\{R,-R,S,-S\}$ in the aligned case and
$\{R,-R,iS,-iS\}$ in the staggered case. In this order the block is
\[
 C=\begin{pmatrix}
 A+B&0&-A&-B\\
 0&-A-B&B&A\\
 -A&B&A-B&0\\
 -B&A&0&B-A
 \end{pmatrix}.
\]
Here aligned rings have
$A=2(R+S)/(R-S)^2$, $B=2(R-S)/(R+S)^2$; staggered rings have
$A=2(R+iS)/(R-iS)^2$, $B=\overline A$.
Let $J$ interchange the first two and last two coordinates. Then
$JCJ=-C$ and $C\mathbf1=0$. The even and odd eigenspaces of $J$
both have dimension two. In an orthonormal real basis of these
spaces, with the first even vector proportional to $\mathbf1$,
$C$ therefore has the form
\[
 \begin{pmatrix}
 0&0&0&0\\0&0&\alpha&\beta\\
 0&\alpha&0&0\\0&\beta&0&0
 \end{pmatrix}.
\]
Its two nonzero singular values are equal, with square
$|\alpha|^2+|\beta|^2=\tfrac12\operatorname{Tr}(CC^*)
=4(|A|^2+|B|^2)$.
For aligned rings, $A=5/7$ and $B^2=8/245$, giving $532/245$.
For staggered rings, $|A|^2=|B|^2=4/(R^2+S^2)$, giving
$32/(R^2+S^2)$. The equal and unequal stationary choices yield
$64/49$ and $8/7$, respectively. Each exceeds one and hence violates
the necessary Hessian contraction proved above.
\end{proof}


% Additive fragment: three entire boundary subcharts of first-e4.
\section{Three exact boundary exclusions in the eight-point first-\texorpdfstring{\(e_4\)}{e4} chart}
\label{sec:first4-three-boundaries}

This fragment proves a new exclusion on three complete subcharts of
\[
 p(X)=X^8+X^4+bX^3+cX^2+dX+u.
\]
The normalization applies when \(e_2=e_3=0\) and the coefficient
\(e_4\ne0\): a nonzero complex scaling makes that coefficient one.
It does not cover arbitrary \(b,c,d\). The conclusion is that a
collision-free physical stationary configuration in any of the three
subcharts
\[
 (b,d)=(0,0),\qquad(c,d)=(0,0),\qquad(b,c)=(0,0)
\]
must have \(b=c=d=0\). Its roots therefore consist of two concentric
squares. The previously proved two-square instability excludes every
such configuration as an unrestricted local or global maximum.

\subsection{Original necessary equations and the reality limitation}
Write \(\delta=\operatorname{Res}_X(p,p')\). Noncollision is exactly
\(\delta\ne0\). At a physical stationary shape of moment \(Q>0\),
\[
 F_i(x)=\sum_{j\ne i}\frac1{x_i-x_j}
       =\frac{28}{Q}\overline{x_i},\qquad
 g_i=\frac{p''(x_i)}{p'(x_i)}=\frac{56}{Q}\overline{x_i}.
\]
Thus \(x_i=2F_i(g)\), and the \(g_i\) are distinct. Define the
original resultant generators by
\[
 R_0(Y)=\operatorname{Res}_X(p,Yp'-p''),\quad
 R_t(Y)=\left.\partial_T\operatorname{Res}_X
          (p,(Y+TX)p'-p'')\right|_{T=0},
\]
\[
 E_1=[Y^5]R_t,\qquad
 E_r=[Y^{6-r}]R_t-(8-r)(7-r)[Y^{8-r}]R_0
       \quad(2\le r\le6).
\]
Indeed, for \(H(Y)=\prod_i(Y-g_i)\), the double-force identity gives
\[
 \frac{R_t(Y)}{\delta}
 =\sum_i\frac{x_iH(Y)}{Y-g_i}
 =2\sum_{i<j}\frac{H(Y)}{(Y-g_i)(Y-g_j)}
 =H''(Y).
\]
Hence \(R_t=R_0''\) and every physical stationary shape has
\(E_1=\cdots=E_6=0\). These equations are only necessary. They do
not impose actual conjugate pairing or make a complex algebraic
solution physical. No sufficiency assertion, numerical coverage, or
claim about positive-dimensional physical branches is used below.
These generators use the double-force first variation, distinct from
the squared-force resultant in the general even-chart model.

\subsection{The full even boundary \texorpdfstring{\(b=d=0\)}{b=d=0}}
Put \(q=c^2\) and define
\begin{align*}
 D={}&-27q^2+144qu-4q+256u^3-128u^2+16u,\\
 A={}&-27q^2-432qu-4q+1792u^3+1664u^2-16u.
\end{align*}
Exact specialization of the original six generators gives
\[
 \delta=256uD^2,\qquad E_1=E_3=E_5=0,\qquad
 E_2=6144cAD,\quad E_4=-8192F,\quad E_6=16384cG.
\]
Here the polynomials \(F,G\), all six specialized original generators,
and the exact division quotients are recorded without omission in
\texttt{even\_original.json} and \texttt{even\_certificate.json}.
Assume \(c\ne0\) and \(\delta\ne0\). Then \(uD\ne0\), and \(E_2=0\)
forces \(A=0\). Exact polynomial division yields
\[
 F=a_FA+\frac{256u}{3}f,\qquad
 G=a_GA-\frac{32}{3}g,
\]
where
\begin{align*}
 f={}&(362880u^4+5241888u^3+3679416u^2+5478u+1)q\\
    &-16128000u^5-27921408u^4-11577856u^3+154816u^2-584u,\\
 g={}&(28340928u^4+33146928u^3-3734172u^2-17271u-28)q\\
    &-169666560u^5-91799040u^4+13605760u^3-134560u^2-112u.
\end{align*}
Consequently \(f=g=0\). For \(M=u(12u+1)^4\), the exact resultants are
\[
 \operatorname{Res}_q(f,g)=24M H_0,\qquad
 \operatorname{Res}_q(A,f)=48M H_1,
\]
with
\begin{align*}
 H_0={}&794740800u^4+545326096u^3-28268724u^2-237885u-686,\\
 H_1={}&237081600u^6-2605639680u^5-4050501504u^4\\
       &-1508424768u^3-25375852u^2-452060u-49.
\end{align*}
The explicit rational polynomials
\(\alpha,\beta\) in \texttt{even\_certificate.json}
(degrees five and three) certify the formal identity
\[
 \alpha H_0+\beta H_1=1,\qquad
 \frac{\alpha}{24}\operatorname{Res}_q(f,g)
 +\frac{\beta}{48}\operatorname{Res}_q(A,f)=u(12u+1)^4.
\]
The coefficient lists, rather than a computed gcd or a Groebner
statement, are the certificate. Since the two resultants vanish and
\(u\ne0\), one has \(u=-1/12\). But
\[
 f(q,-1/12)=\frac{7360}{9}(27q+8),
\]
so \(q=-8/27\), and \(D(-8/27,-1/12)=0\), contradicting
\(\delta\ne0\). Thus \(c=0\).

\subsection{The full \texorpdfstring{\(b\)}{b}-only boundary \texorpdfstring{\(c=d=0\)}{c=d=0}}
Assume \(b\ne0\), and put \(v=b^4\). Here
\[
 E_1=4096b^3u^2(20u+27)(60u+1)^2,\qquad u^2\mid\delta.
\]
Noncollision implies \(u\ne0\); hence \(u=-27/20\) or \(-1/60\).
The relevant exact generators are
\[
 E_2=1536b^2u^2B,\qquad E_5=2304b^3uC,
\]
\begin{align*}
 B={}&(-112500u+13125)v+
       2662400u^3-276480u^2-98816u+384,\\
 C={}&(1687500u-16875)v+
       232243200u^3+2498560u^2+326656u.
\end{align*}
All divided factors are nonzero under the stated assumptions, so
\(B=C=0\). At \(u=-27/20\), their exact values satisfy
\[
 B=165000v-\frac{34603008}{5},\quad
 C=-2295000v-\frac{2836463616}{5},\quad
 -\frac{17B}{811008000}-\frac{C}{663552000}=1,
\]
which is impossible. At \(u=-1/60\),
\[
 B=15000v+\frac{262144}{135},\quad
 C=-45000v-\frac{262144}{45},
\]
so \(v=-32768/253125\). Direct exact substitution in the original
discriminant gives
\[
 \delta(-1/60,-32768/253125)=0.
\]
This is the second contradiction.
Thus \(b=0\). Full specialized generator and discriminant formulas are
in \texttt{b\_only\_original.json}; the two branch certificates are in
\texttt{single\_odd\_reduce.json}.

\subsection{The full \texorpdfstring{\(d\)}{d}-only boundary \texorpdfstring{\(b=c=0\)}{b=c=0}}
Assume \(d\ne0\), and put \(v=d^4\). Three exact generators suffice:
\[
 E_2=-1536d^2B,\qquad E_3=-1024d^3C,\qquad
 E_6=-917504d^2H,
\]
\begin{align*}
 B={}&(-470596u+50421)v+
       2408448u^4-716800u^3+37888u^2-2304u,\\
 C={}&2470629v+6146560u^4-11283328u^3
           +1378272u^2-219240u+10368,\\
 H={}&(-941192u-151263)v+5531904u^4+344960u^3-67440u^2.
\end{align*}
Their nonzero scalar factors can be divided out, hence \(B=C=H=0\).
For \(N=(28u+9)^2\), the two exact resultants are
\[
 \operatorname{Res}_v(B,C)=-134456N K_0,\qquad
 \operatorname{Res}_v(C,H)=134456N K_1,
\]
\[
 K_0=27440u^3-14504u^2+1239u-48,\qquad
 K_1=54880u^3+2450u^2-3045u+144.
\]
The following explicit rational polynomials certify
\(\alpha_dK_0+\beta_dK_1=1\):
\begin{align*}
 \alpha_d={}&\frac{89622225}{80734112}u^2+
             \frac{358103305}{3875237376}u-
             \frac{469880965}{7750474752},\\
 \beta_d={}&-\frac{89622225}{161468224}u^2+
             \frac{263473525}{968809344}u-
             \frac{102804247}{7750474752}.
\end{align*}
Therefore the two vanishing resultants force \((28u+9)^2=0\), or
\(u=-9/28\). At this value,
\[
 C=3(823543v+221184),\qquad
 \delta=-\frac{(823543v+221184)^2}{823543}.
\]
Thus \(C=0\) contradicts \(\delta\ne0\). This proves \(d=0\),
including the possible origin boundary \(u=0\), without dividing by
\(u\). Full original formulas are in \texttt{d\_only\_original.json}.

\subsection{Consequence and exact replay}
In all three subcharts the only noncollision physical stationary
possibility is \(p=X^8+X^4+u\). Its two distinct nonzero fourth-power
roots give two concentric square rings. The existing physical
two-square theorem, proved in
\texttt{source/n8/n8\_next\_even\_model.tex},
gives symmetry-breaking singular-value squares
\(64/49\), \(532/245\), or \(8/7\), each greater than one.
Hence no shape in these three whole subcharts is an unrestricted
local maximum.

\texttt{reconstruct\_chart.py} reconstructs each subchart from \(p\)
using the exact \(15\)-by-\(15\) Sylvester determinant. The derivative
is the sum of eight row-replacement determinants with
\(q_1=Xp'-8p\); this is equivalent modulo \(p\) to \(Xp'\).
Each subchart passes ten formal identities, including comparison of
all six \(E_r\) and \(\delta\) with the saved original first-\(e_4\)
generators. No points are sampled.
\texttt{verify\_certificates.py} uses only standard-library sparse
polynomial arithmetic over \(\mathbb Q\) to replay \(46\) formal
factorizations, division identities, resultant formulas and Bezout
certificates. It passes both normally and with \texttt{-S -O}.
The unrestricted first-\(e_4\) chart with multiple nonzero coefficients
\(b,c,d\) remains unresolved by this fragment.

\section{An exact phase reduction for the unresolved \texorpdfstring{\(1+3+4\)}{1+3+4} branch}
\label{sec:134-generic-phase}

This section supplies necessary equations for a remaining branch; it
does not exhaust it. Use coordinates \(u\) with moment seven and force
constant four. Suppose a no-origin physical stationary configuration
has exactly three distinct squared radii \(r,c,e\), with populations
\(1,3,4\), and all three radial group sums are nonzero. The two sum
identities \(\sum u_i=\sum |u_i|^2u_i=0\) make these sums collinear.
Rotate their common line so that the singleton is \(x>0\), and put
\[
 r=x^2=7-3c-4e,\quad
 a=x\frac{e-r}{c-e},\qquad b=x\frac{r-c}{c-e}.
\]
The triangle and quadrilateral root polynomials are necessarily
\[
 T(X)=X^3-aX^2-\frac{af}{c}X+f,\qquad |f|^2=c^3,
\]
\[
 G(X)=X^4-bX^3+qX^2-\frac{bd}{e}X+d,
 \qquad |d|^2=e^4,\quad e^2q=d\overline q.
\]
For actual global candidates the existing radius bound gives
\(0<c,e<49/32<13/8\).

For \(p_j=\sum_i u_i^j\), stationarity gives, by pairwise summation,
\begin{equation}\label{eq:134-new-recursion}
 \sum_i |u_i|^2u_i^j
 =\frac{15-j}{8}p_j+
       \frac18\sum_{h=1}^{j-1}p_hp_{j-h}\qquad(j\ge1).
\end{equation}
The imaginary parts for \(j=2,3\) are
\[
 (c-13/8)\frac ac\Im f-(e-13/8)\Im q=0,
\]
\[
 (c-3/2)(a^2/c-1)\Im f+
             (e-3/2)b(-\Im q+\Im d/e)=0.
\]
In particular \(\Im f\ne0\) forces \(\Im q\ne0\).

\begin{lemma}[The entire degenerate triangle branch is reflected]
\label{lem:134-new-degenerate-reflected}
For a physical stationary \(1+3+4\) configuration as above, if
\(a^2=c\), then the entire configuration is invariant under conjugation.
This assertion does not require the global radius bound.
\end{lemma}
\begin{proof}
The triangle factors as \(T=(X-a)(X^2-f/a)\). Centering gives
\(b=-x-a\ne0\), and the first radial sum gives
\[
 e=\frac{x^3+a^3}{x+a}=x^2-xa+a^2.
\]
Put \(t=x/a\). The moment constraint then gives
\[
 c=\frac7{5t^2-4t+7}\le\frac{35}{31}<\frac{13}{8},\qquad
 e=\frac{7(t^2-t+1)}{5t^2-4t+7},
\]
\[
 \frac32-e=\frac{c}{14}\bigl((t+1)^2+6\bigr)>0.
\]
Thus both low imaginary identities have nonzero nonresonant factors.
If \(\Im f=0\), the second moment gives \(\Im q=0\), and the third
gives \(\Im d=0\), since \(a^2/c-1=0\) and \(b\ne0\).
All coefficients are then real and the configuration is reflected.

Suppose instead \(\Im f\ne0\); then \(\Im q\ne0\). The third
moment gives \(\Im d=e\Im q\). The circle identity
\(e^2q=d\overline q\) now forces
\[
 d=e(q-e),\qquad
 G=(X^2-bX+e)(X^2+q-e).
\]
Indeed its imaginary part gives \(\Re d=e\Re q-e^2\).
The triangle and quadrilateral each consist of one antipodal pair
together with a reflected remainder. The other four roots have
polynomial \((X-x)(X-a)(X^2-bX+e)\) with real coefficients.
Consequently all odd global power sums and all their radial weighted
sums are real. In particular
\[
 p_3=x^3+a^3+b^3-3eb=3(x+a)(x-a)^2\ne0.
\]
Here \(x+a\ne0\) because \(b\ne0\), and \(x-a\ne0\) because
\(r\ne c\). The fifth moment identity has imaginary part
\(p_3\Im p_2/4=0\), so \(\Im p_2=0\).
But Newton's identity gives
\[
 \Im p_2=2\left(\frac{\Im f}{a}-\Im q\right),\qquad
 \Im q=\frac1a\frac{c-13/8}{e-13/8}\Im f.
\]
These imply \(c=e\), contradicting the distinct radial levels.
Thus the nonreflected possibility is impossible.
\end{proof}

\begin{proposition}[Generic phases determined by the two radii]
\label{prop:134-new-phases}
Within this global radius range, every nonreflected physical
configuration in this branch has \(\Im f\ne0\), \(a^2\ne c\) and
\(e\ne3/2\). Its real parts
\(\Re f/x\) and \(\Re q\) are uniquely determined rational functions
of \(c,e\) by the fourth and fifth moment identities.
Thus the generic phase reduction covers every nonreflected candidate
in the stated radius range, although it does not exclude them.
\end{proposition}
\begin{proof}
The lemma rules out \(a^2=c\) for a nonreflected configuration.
We first show that it has \(\Im f\ne0\). If \(\Im f=0\), then
\(\Im q=0\). For \(e\ne3/2\), the third imaginary identity gives
\(\Im d=0\), so it is reflected. For \(e=3/2\) and \(q\ne0\),
the circle identity also gives \(\Im d=0\). It remains to treat
\(e=3/2,q=0\). If \(\Im d\ne0\), the global second power sum is
real and the fourth moment identity forces the quadrilateral fourth
power sum \(b^4+4d(b^2/e-1)\) real. Thus \(b^2=e\).
But \(r+3c=1\), \(0<c<1/3\), and the radial sum formula gives
\[
 |b|=\frac{\sqrt{1-3c}\,|1-4c|}{3/2-c}
       <\frac67<\sqrt{3/2},
\]
since \(|1-4c|\le1\) and \(3/2-c>7/6\).
This contradicts \(b^2=e\). Hence this resonance is reflected too,
and a nonreflected configuration must have \(\Im f\ne0\).

Now if \(e=3/2\), the third imaginary identity forces \(a^2=c\):
the moment constraint gives \(c<1/3\), hence \(c\ne3/2\).
The lemma excludes this. All three stated nondegeneracy conditions
therefore hold. Write
\[
 \alpha=\frac{e-r}{c-e},\quad \beta=\frac{r-c}{c-e},\quad
 \kappa=\frac{\alpha(c-13/8)}{c(e-13/8)},
\]
\[
 \lambda=e\left[\kappa-
    \frac{(c-3/2)(r\alpha^2/c-1)}{(e-3/2)r\beta}\right].
\]
All displayed denominators are nonzero: \(r,c,e\) are positive and
distinct, and \(a,b\ne0\). Put \(h=\Im f\ne0\),
\(U=\Re f/x\), \(Q=\Re q\). The two low imaginary identities and
the imaginary part of \(e^2q=d\overline q\) give
\begin{equation}\label{eq:134-new-parameters}
 f=xU+\mathrm i h,\quad q=Q+\mathrm i x\kappa h,\quad
 d=\frac{\lambda}{\kappa}Q-e^2+\mathrm i x\lambda h.
\end{equation}

Let \(t_j,g_j\) be the power sums of the monic polynomials \(T,G\),
respectively, computed by Newton's identities, and set
\(p_j=x^j+t_j+g_j\). Define the explicit polynomial expressions
\[
 Z_j=x^{j+2}+ct_j+eg_j-\frac{15-j}{8}p_j
                  -\frac18\sum_{h'=1}^{j-1}p_{h'}p_{j-h'}.
\]
Substitution of \eqref{eq:134-new-parameters} and \(x^2=r\) gives
\[
 \begin{pmatrix}\Im Z_4/(xh)\\\Im Z_5/h\end{pmatrix}
       =M(c,e)\begin{pmatrix}U\\Q\end{pmatrix}+N(c,e).
\]
Both entries are affine in \(U,Q\), and independent of \(h\).
This can also be seen directly from Newton's identities: through
degree five the triangle sums contain only \(f,f^2\), and the
quadrilateral sums only \(q,q^2,d,dq\); their imaginary parts after
division by \(h\) are affine. The moment products on the right have
the same property because \(p_1=0\).
Here \(N\) is unambiguously obtained from the displayed finite
recurrence by setting \(U=Q=0\); no solving or root sampling defines
it. For compact explicit matrix formulas put
\[
 K=7c^2+24ce-35c+25e^2-70e+49,
\]
\[
 L_1=80ce-126c-124e+195,\qquad
 L_2=80ce-124c-126e+195.
\]
Direct expansion gives
\begin{align*}
 M_{11}&=\frac{r(8c-13)(8e-11)(e-r)^2}
                    {2c^2(c-e)^2(8e-13)},&
 M_{12}&=\frac{(8c-11)(e-r)}{2c(c-e)},\\
 M_{21}&=\frac{r(e-r)(r-c)(4e-5)L_1K}
                  {2c^2(c-e)^3(2e-3)(8e-13)},&
 M_{22}&=\frac{(4c-5)(r-c)L_2K}
                  {2c(c-e)^2(8c-13)(2e-3)}.
\end{align*}
Their exact determinant factors as
\begin{equation}\label{eq:134-new-determinant}
 \det M=-\frac{r(e-r)^2(r-c)JK}
               {2c^3(c-e)^3(2e-3)(8e-13)},\qquad
 J=128ce-208c-208e+335.
\end{equation}
One has the strictly positive global-domain bound
\[
 J\ge\frac{297}{32}>0
 \quad\text{on}\quad
 0\le c,e\le49/32,\quad 3c+4e\le7.
\]
Indeed \(J\) decreases in each variable on this polygon. On its
oblique moment edge, substitution \(e=(7-3c)/4\) makes \(J\) a
concave quadratic, whose minimum is at an endpoint. The five vertex
values are \(335,33/2,297/32,13,33/2\).
Finally the exact identity
\[
 (c-e)^2(a^2-c)=(r-c)K
\]
shows that \(K\ne0\) on the generic branch. Thus \(\det M\ne0\),
and \(M(U,Q)^{\mathsf T}+N=0\) determines both real parts rationally.
This proves the proposition.
\end{proof}

The remaining necessary circle constraints now include the explicit
two-radius conditions
\[
 h^2=c^3-rU^2>0,\qquad
 \frac{\lambda}{\kappa}
       \bigl(Q^2+r\kappa^2h^2\bigr)-2e^2Q=0.
\]
The second is the real part of \(e^2q=d\overline q\); together with
\eqref{eq:134-new-parameters} it also gives \(|d|^2=e^4\).
The signs \(h\) and \(-h\) give conjugate configurations. These
necessary conditions still require the full force remainders and
actual common-circle root tests; they are not sufficient physical
equations and do not establish finite coverage.
In the exceptional branch, \(a^2=c\) factors the triangle explicitly:
\[
                  T(X)=(X-a)(X^2-f/a).
\]
It consists of a real point and an antipodal pair on that circle.
The lemma eliminates its entire nonreflected part. A value or force
exhaustion for reflected configurations of this population, and the
remaining generic two-radius force problem, are still open here.

The independently reconstructed checker
\path{phase_reduction_check.py} uses exact SymPy arithmetic for ten
formal tests: the original Newton recurrences, affine dependence,
the determinant factorization, the degeneracy identity, and all
vertex signs. No numerical roots or sampling certify this reduction.
The independent companion-matrix program \path{phase_independent_audit.py}
also reconstructs the low force moments, checks every printed matrix
entry, both reflected degeneracy arguments, and the additional
radius curve below. Its 57 exact checks pass independently of the
Newton-recurrence implementation. A further 22-identity audit supplies
a Bernstein certificate for the whole resonant radius interval.

% Appendable supplement after the generic I4/I5 phase reduction.
\begin{proposition}[A degree-six curve for the remaining radial parameters]
\label{prop:eight-134-new-real-radius-curve}
In the nonreflected physical $1+3+4$ branch in the stated global
radius range, put $r=7-3c-4e$, $U=\Re f/x$, $Q=\Re q$, and retain
the established parameters $\kappa,\lambda$.
Define
\begin{align*}
 P_2={}&24c^3+192c^2e-309c^2+248ce^2-950ce+833c\\
       &\quad+48e^3-469e^2+1015e-637,\\
 P_3={}&36c^3+162c^2e-213c^2+236ce^2-630ce+406c\\
       &\quad+126e^3-487e^2+609e-245,\\
 K={}&7c^2+24ce-35c+25e^2-70e+49,\\
 P_6={}&3(2c-3)KP_2-(8c-13)(e-r)(c-e)P_3.
\end{align*}
Every such physical stationary configuration satisfies
\[
                         P_6(c,e)=0.
\]
This is an additional necessary radius condition, not an exhaustion
of its real roots or of physical stationary configurations.
\end{proposition}
\begin{proof}
Let $Z_j$ denote the mixed force remainder already defined in the
phase reduction. Under $f=xU+\mathrm ih$, $q=Q+\mathrm ix\kappa h$,
$d=(\lambda/\kappa)Q-e^2+\mathrm ix\lambda h$, the two low real
moments have the form
\[
 \Re Z_2=A_2U+B_2Q+N_2,\qquad
 \Re Z_3/x=A_3U+B_3Q+N_3.
\]
Exact expansion, or companion-matrix traces, gives
\begin{align*}
 A_2&=\frac{r(e-r)(8c-13)}{4c(c-e)},&
 B_2&=-\frac{8e-13}{4},&
 N_2&=\frac{rP_2}{4(c-e)^2},\\
 A_3&=\frac{3(2c-3)(r-c)K}{2c(c-e)^2},&&&
 N_3&=\frac{(r-c)P_3}{2(c-e)^2},\\
 B_3&=-\frac{3(2c-3)(8e-13)(r-c)K}
               {2(c-e)(8c-13)r(e-r)}.
\end{align*}
The two coefficient rows have the same homogeneous null direction:
\[
 A_2B_3-A_3B_2=0.
\]
Their constant compatibility is exactly
\[
 A_2N_3-A_3N_2
       =-\frac{r(r-c)}{8c(c-e)^4}\,P_6(c,e).
\]
The force equations set both real remainders to zero. Since
$r>0$, $c>0$, and the three radial levels are distinct, the
displayed nonzero prefactor forces $P_6=0$.
Every identity here was independently reconstructed from companion
traces in \path{phase_independent_audit.py}; the checker also verifies
all six low-real coefficients separately.
\end{proof}

For possible further elimination, the second real moment and fourth
imaginary moment give another invertible pair. Its exact determinant is
\[
 \det\begin{pmatrix}A_2&B_2\\M_{11}&M_{12}\end{pmatrix}
       =\frac{r\alpha^2(8c-13)(4c+4e-11)}{4c^2}.
\]
On the global polygon,
$c+e\le7/4+c/4\le273/128<11/4$, so this determinant does not
vanish. This observation is recorded as a further algebraic route;
it does not replace the established fourth/fifth imaginary-moment
parameterization or prove that the remaining radius curve is empty.

\section{What is now complete and what remains}
With the stated published classification, the six- and seven-point
unrestricted problems, including all no-origin configurations, are
complete. The classification-free replacements are separate research
questions; their incompleteness does not reopen these global values.

Eight points are still incomplete. A remaining no-origin global
candidate has at least three radial levels. If centrally symmetric,
it has four distinct levels and no reflection axis. Three-radius
configurations with nonzero collinear radial sums can still have
populations \(1+2+5\), \(1+3+4\), \(2+2+4\), or \(2+3+3\).
The \(1+3+4\) reduction above excludes nonreflected degeneracy, but
does not exhaust its remaining two-radius force system or its
reflected configurations.
General multiple-coefficient first-fourth charts and configurations
with four or more radial levels are also not exhausted.

The earlier fixed-origin theorem proves its entire eight-point
comparison. The companion regional addendum excludes the explicit
near-origin region for actual global maximizers. These narrow the remaining
no-origin problem. No finite sampling, necessary algebraic equation,
or successful inequality check is substituted for missing coverage.

\begin{thebibliography}{9}
\bibitem{FS}
J.-C.~Faug\`ere and J.~Svartz,
\emph{Solving polynomial systems globally invariant under an action of
the symmetric group and application to the equilibria of \(N\) vortices
in the plane}, Proceedings of ISSAC 2012, 170--178.
\href{https://doi.org/10.1145/2442829.2442856}{doi:10.1145/2442829.2442856}.
\bibitem{FSdata}
J.-C.~Faug\`ere and J.~Svartz, author computational archive,
\emph{Vortices for \(N=6\)}, exact coefficient bases supplied as
\path{archivepolsys_vortices_N6.txt}.
\bibitem{FSdataSeven}
J.-C.~Faug\`ere and J.~Svartz, author computational archive,
\emph{Vortices for \(N=7\)}, retained exact coefficient bases supplied as
\path{archivepolsys_vortices_N7.txt}.
\end{thebibliography}
\end{document}
